% ==================================================================
% id:       lim-020
% topic:    גבולות ורציפות של פונקציות
% subtopic: רציפות ונקודות אי־רציפות
% solution: full
% ==================================================================
\textbf{
הוכיחו או הפריכו
כל אחד מהסעיפים הבאים:}
\begin{enumerate}
    \item אם
    $f$
    רציפה
    בקטע
    $I$,
    אז
    $f$
    חסומה ב-
    $I$.
\hint{בדקו את הטענה עבור סוגים שונים של קטעים (סגור, פתוח, לא חסום); אם נראה לכם שהיא לא נכונה, חפשו דוגמה נגדית פשוטה.}
\begin{solbox}
אם
$I$
הוא קטע סגור, אז הטענה נכונה
(משפט וויירשטראס),
אך
$I$
לא בהכרח סגור. נספק דוגמא נגדית.

ניקח
$f(x)=\tan x$
בקטע
$I=(-\tfrac{\pi}{2},\tfrac{\pi}{2})$.
הפונקציה רציפה בקטע זה אך בקצה הימני שואפת לאינסוף ובקצה השמאלי למינוס אינסוף,
לכן אינה חסומה. לסיכום,
הטענה אינה נכונה.
\end{solbox}
    \item
    אם
    $f$
    פונקציה לא רציפה
    בקטע
    $[a,b]$, אז
    $f$
    לא חסומה בקטע $[a,b]$.
\hint{חשבו על פונקציות לא רציפות שאתם מכירים ובדקו אם הן חסומות; אם הטענה נכונה, נסו להוכיח אותה.}
\begin{solbox}
הטענה אינה נכונה.

ניקח
$f(x)=D(x)$ -
פונקציית דיריכלה
בקטע
$[0,1]$.
זאת פונקציה לא רציפה
(כל הנקודות בקטע הינן נקודות אי-רציפות שלה),
אבל
מתקיים
$$,0\leq D(x)\leq 1$$
כלומר,
$D(x)$
חסומה ב-
$[0,1]$.
\end{solbox}
    \item
    אם
    $f$
    רציפה ב-$[a,b]$, אז
    $f$
    חסומה בקטע $(a,b)$.
\hint{מה ידוע על פונקציה רציפה בקטע סגור, ומה הקשר בין $[a,b]$ ל-$(a,b)$? אם נראה לכם שהטענה לא נכונה, חפשו דוגמה נגדית פשוטה.}
\begin{solbox}
הטענה נכונה.

ממשפט וויירשטרס נסיק כי
$f$
חסומה בקטע
$[a,b]$,
לכן בפרט חסומה בקטע
$(a,b)$.
\end{solbox}
    \item
    אם
    $f$
    רציפה ב-$[0,1]$,
    $f(0)=2$ ו- $f(1)=5$ אז התמונה של $f$ היא $[2,5]$.
\hint{הפרידו בין שתי ההכלות: האם $[2,5]$ מוכל בתמונה, והאם התמונה מוכלת ב-$[2,5]$? אם אחת מהן נכשלת, חפשו דוגמה נגדית פשוטה.}
\begin{solbox}
הטענה אינה נכונה.

הרעיון הוא לקחת פונקציה עם הערכים הנתונים בקצוות ולגרום לה להשתולל בין לבין, למשל לעלות לערכים גבוהים מ-$5$
או קטנים מ-$2$.
למשל,
ניקח
$f(x)=10x^2-7x+2$.
בבירור
$f$
רציפה ב-
$[0,1]$
ומתקיים
$$,f(0)=2,\ f(1)=5$$
מכיוון ש-
$f(\tfrac{1}{2})=1\notin[2,5]$,
נסיק כי תמונתה של
$f$
אינה הקטע
$[2,5]$.
\end{solbox}
    \item
    אם
    $f$
    רציפה ומונוטונית ב-$[0,1]$, $f(0)=2$ ו- $f(1)=5$ אז התמונה של $f$ היא $[2,5]$.
\hint{בדקו שוב את שתי ההכלות בין התמונה ל-$[2,5]$, והפעם שימו לב מה המונוטוניות מוסיפה.}
\begin{solbox}
הטענה נכונה.

מכיוון ש-
$f$
מונוטונית ומתקיים
$f(0)=2<f(1)=5$,
היא מונוטונית עולה. לכן לכל
$x\in[0,1]$
מתקיים
$$,2=f(0)\leq f(x)\leq f(1)=5$$
ולכן
${\rm{Im}}(f)\subseteq[2,5]$.
מצד שני, עבור
$2=f(0)< y< f(1)=5$
קיים
$c\in[0,1]$
כך ש-
$f(c)=y$
(ממשפט ערך הביניים),
לכן
${\rm{Im}}(f)=[2,5]$.
\end{solbox}
    \item
    אם
    $f$
    רציפה ב-$[0,1]$,
    $f(0)=2$, $f(1)=5$ והתמונה של $f$ היא $[2,5]$ אז $f$ מונוטונית.
\hint{נסו לצייר גרפים שונים שעומדים בכל הדרישות, ובדקו אם כולם חייבים להיות מונוטוניים.}
\begin{solbox}
הטענה אינה נכונה.
ציירו פונקציה רציפה העונה על הדרישות והרסו את המונוטוניות שלה בפנים הקטע, רק אל ״תצאו מהקווים״.

להלן דוגמא נגדית:
$$.f(x)=2+3|\sin(\tfrac{3\pi x}{2})|$$
בבירור
$f$
היא פונקציה רציפה.
נשים לב כי
$$,2\leq f(x)\leq 5$$
בנוסף
$$,f(0)=2,\ f(1)=5$$
כלומר,
$f$
עומדת בדרישות של השאלה. אבל
$$,f(\tfrac{1}{3})=5,\ f(\tfrac{2}{3})=2$$
ולכן
$f$
אינה מונוטונית.
\end{solbox}
    \item
    אם
    $f$
    חסומה
    ב-$[a,b]$,
    אז
    $f$
    רציפה בקטע $[a,b]$.
\hint{השוו לסעיף ב׳, שבו הכיוון הפוך. אם נראה לכם שהטענה לא נכונה, חפשו דוגמה נגדית פשוטה; אחרת נסו להוכיח אותה.}
\begin{solbox}
הטענה אינה נכונה.
ניקח שוב את הדוגמא של סעיף ב'.
\end{solbox}
    \item
    אם
    $f$
    חסומה
    ב-$[a,b]$,
    אז
    $f$
    מקבלת מקסימום בקטע $[a,b]$.
\hint{בדקו אם חסימות לבדה מבטיחה שהחסם העליון מושג. אם נראה לכם שהטענה לא נכונה, חפשו דוגמה נגדית פשוטה; אחרת נסו להוכיח אותה.}
\begin{solbox}
הטענה אינה נכונה.
תהא
$$.\ds{f(x)=\case{cc}{x & x\in[0,1) \\ -3 & x=1}}$$
הפונקציה חסומה בקטע
$[0,1]$,
אך אינה מקבלת ערך מקסימלי.
אכן,
לכל
$x\in[0,1)$
$f(x)>-3=f(1)$
ולכל
$0<x_0\in[0,1)$
מתקיים
$\frac{x_0+1}{2}<1$
ו-
$$,f(\tfrac{x_0+1}{2})=\frac{x_0+1}{2}>x_0=f(x_0)$$
כלומר, לכל ערך נתון שהפונקציה מקבלת בקטע הנדון, ניתן לספק
$x$
מהקטע כך שערכה של
$f$
ב-
$x$
יהיה גדול יותר,
זאת אומרת שאין ל-
$f$
מקסימום בקטע
$[0,1]$.
\end{solbox}
\end{enumerate}
